\section*{第 \ref{ch:b3:many-electron-atoms} 章——多电子原子与谱项符号}

\begin{solution}{exo:b3:many-electron-atoms:1}
\[
\Psi = \frac{1}{\sqrt6}\begin{vmatrix}1s\alpha(1) & 1s\beta(1) & 2s\alpha(1)\\
1s\alpha(2) & 1s\beta(2) & 2s\alpha(2)\\ 1s\alpha(3) & 1s\beta(3) & 2s\alpha(3)
\end{vmatrix}.
\]
$1s$ 只有两个自旋轨道 $1s\alpha$ 和 $1s\beta$，第三个电子
会重复其中之一，使两列相等，行列式为零。
\end{solution}

\begin{solution}{exo:b3:many-electron-atoms:2}
$\binom62 = 15$，$\binom63 = 20$，$\binom{10}{2} = 45$。$p^3$：${}^4$S（4）+
${}^2$D（10）+ ${}^2$P（6）= 20。$d^2$：${}^3$F（21）+ ${}^3$P（9）+ ${}^1$G（9）+
${}^1$D（5）+ ${}^1$S（1）= 45。
\end{solution}

\begin{solution}{exo:b3:many-electron-atoms:3}
N（$2p^3$，半满）：\termsym{4}{S}{3/2}。O（$2p^4$，超过半满）：
\termsym{3}{P}{2}。F（$2p^5$）：\termsym{2}{P}{3/2}。\ce{Ti^2+}（$3d^2$）：
\termsym{3}{F}{2}。\ce{Fe^3+}（$3d^5$，半满，自旋全部平行，$L = 0$）：
\termsym{6}{S}{5/2}.
\end{solution}

\begin{solution}{exo:b3:many-electron-atoms:4}
${}^2$D：$J = \frac52$（6 个状态）和 $\frac32$（4 个状态）；$6 + 4 = 10 = 5 \times 2$。
${}^3$P：$J = 2$（5），1（3），0（1）；$5 + 3 + 1 = 9 = 3 \times 3$。
\end{solution}

\begin{solution}{exo:b3:many-electron-atoms:5}
能级间隔 $16.4$（$J = 1 - 0$）和 $27.0$（$2 - 1$）：比值为 1.65，而非 2。
由第一个得 $A = 16.4/1 = \qty{16.4}{cm^{-1}}$，由第二个得 $27.0/2 = \qty{13.5}{cm^{-1}}$。
单一的 $A$ 无法拟合：这些能级与其他谱项
（${}^1$D、${}^1$S）的能级略有混合，而罗素–桑德斯耦合只是一种
近似。
\end{solution}

\begin{solution}{exo:b3:many-electron-atoms:6}
$\lambda = 10^7/\tilde\nu$：\qty{589.76}{nm}（D$_1$）和 \qty{589.16}{nm}
（D$_2$）。间隔 $\qty{17.20}{cm^{-1}} = \qty{2.13}{meV}$。$E(\frac32) -
E(\frac12) = \frac32A$，因此 $A = \qty{11.47}{cm^{-1}}$。
\end{solution}

\begin{solution}{exo:b3:many-electron-atoms:7}
$3p \to 3s$：允许（$\Delta l = -1$）。$3d \to 3s$：禁戒（$\Delta l = -2$）。
$3d \to 3p$：允许。$4s \to 3s$：禁戒（$\Delta l = 0$，
宇称不变）。$4s \to 3p$：允许。
\end{solution}

\begin{solution}{exo:b3:many-electron-atoms:8}
${}^3$P 的平均值：$(5 \times 169086.76 + 3 \times 169086.84 + 169087.83)/9 =
\qty{169086.91}{cm^{-1}}$。与 ${}^1$P 的间隔：\qty{2047.98}{cm^{-1}} $=
\qty{0.254}{eV}$，故 $K(1s,2p) = \qty{0.127}{eV}$，
仅为 $K(1s,2s)$ 的三分之一。交换积分度量的是两个轨道
电子密度的重叠程度：$2s$ 轨道会穿入 $1s$ 区域（在原子核附近
有电子分布），$2p$ 轨道在原子核处为零，重叠较少。
\end{solution}

\begin{solution}{exo:b3:many-electron-atoms:9}
最大的 $M_L = 4$ 要求两个电子都处于 $m_l = 2$，且自旋相反，
得单重态 ${}^1$G（9 个状态）。下一个 $M_L = 3$ 由 $m_l = 2, 1$ 给出，
可取 $M_S = 1$：${}^3$F（21 个状态）。剩下的 $M_L = 2$ 只与 $M_S = 0$ 配对：
${}^1$D（5 个状态）。剩下的 $M_L = 1$ 取 $M_S = 1$：${}^3$P（9 个状态）。还剩一个
状态，$M_L = M_S = 0$，即 ${}^1$S。总计 45。
\end{solution}

\begin{solution}{exo:b3:many-electron-atoms:10}
当 $\mathbf r_1 = \mathbf r_2 = \mathbf r$ 时，三重态波函数为
$\frac{1}{\sqrt2}[1s(\mathbf r)2s(\mathbf r) - 2s(\mathbf r)1s(\mathbf r)] = 0$；
而单重态波函数为 $\sqrt2\,1s(\mathbf r)2s(\mathbf r) \ne 0$。在三重态中，
两个电子永远不会出现在同一点，平均而言
相距更远，因此排斥作用较小：$E_+ - E_- = 2K > 0$。
\end{solution}

\begin{solution}{exo:b3:many-electron-atoms:11}
$d^8$ 超过半满，耦合关系倒转，$J = L + S = 4$
的能级最低。能级间隔：$E(3) - E(4) = 1360.7$ 与 $E(2) - E(3) = 908.9$；比值
1.50，而朗德间隔定则给出 $4/3 = 1.33$。$|A| = 1360.7/4 = \qty{340}{cm^{-1}}$
（$A < 0$）。在这种较重的离子中，自旋–轨道耦合远大于
碳原子的情况（\qty{16}{cm^{-1}}）。
\end{solution}

\begin{solution}{exo:b3:many-electron-atoms:12}
$\sum_J(2J+1)[J(J+1) - L(L+1) - S(S+1)] = 0$（$(2L+1)(2S+1)$ 个状态可以
按任一基底计数，并且在未耦合的基底上 $\sum M_LM_S = 0$）。
对于 ${}^3$P：$1(0 - 4) + 3(2 - 4) + 5(6 - 4) = -4 - 6 + 10 = 0$。碳的
加权平均为 $(0 + 3 \times 16.42 + 5 \times 43.41)/9 = \qty{29.6}{cm^{-1}}$，
即在无自旋–轨道耦合时该谱项的能量。
\end{solution}

\begin{solution}{pb:b3:many-electron-atoms:1}
\textbf{1.} $1s\alpha$, $1s\beta$, $2s\alpha$, $2s\beta$.
\textbf{2.} $|1s\alpha\,2s\alpha|$, $|1s\alpha\,2s\beta|$, $|1s\beta\,2s\alpha|$,
$|1s\beta\,2s\beta|$.
\textbf{3.} $|1s\alpha\,2s\alpha| = \frac{1}{\sqrt2}[1s(1)2s(2) - 2s(1)1s(2)]\,
\alpha(1)\alpha(2)$，将 $2\times2$ 行列式展开即得。
\textbf{4.} 二者之和为 $\frac{1}{\sqrt2}[1s(1)2s(2) - 2s(1)1s(2)]\cdot
\frac{1}{\sqrt2}[\alpha(1)\beta(2) + \beta(1)\alpha(2)]$（乘以 $\sqrt2$，再
归一化）；二者之差为 $\frac{1}{\sqrt2}[1s(1)2s(2) +
2s(1)1s(2)]\cdot\frac{1}{\sqrt2}[\alpha(1)\beta(2) - \beta(1)\alpha(2)]$.
\textbf{5.} 空间部分对称、自旋函数反对称（$S = 0$）；
空间部分反对称，配以三个对称的自旋函数（$S = 1$）。
\textbf{6.} $2S + 1 = 1$ 与 3，即能量相同的状态分别为一个和三个。
\textbf{7.} $[\ce{Ar}]\,3d^2$；$\binom{10}{2} = 45$ 个行列式。
\textbf{8.} 最大 $S = 1$，最大 $L = 2 + 1 = 3$，未过半满：
\termsym{3}{F}{2}，即数据中位于 0 处的能级。
\textbf{9.} $21 + 9 + 9 + 5 + 1 = 45$.
\textbf{10.} 正常次序（$J = 2$ 最低）。能级间隔 184.9 与 235.5：比值 1.27
与 $4/3 = 1.33$ 相比；该定则成立，误差在 5\,\% 以内。
\textbf{11.} ${}^3$F: $(5 \times 0 + 7 \times 184.9 + 9 \times 420.4)/21 =
\qty{241.8}{cm^{-1}}$; ${}^3$P: $(10538.4 + 3 \times 10603.6 + 5 \times
10721.2)/9 = \qty{10661.7}{cm^{-1}}$.
\textbf{12.} $15B = \qty{10419.9}{cm^{-1}}$, $B = \qty{695}{cm^{-1}}$.
\textbf{13.} $\Delta S = 0$.
\textbf{14.} $\Delta S = 1$ 禁戒，$1s2s \to 1s^2$ 中 $\Delta l = 0$
（宇称不变；${}^3$S$_1 \to {}^1$S$_0$ 在两个 S
态之间同样有 $\Delta L = 0$）。该态为亚稳态，其寿命
远长于普通激发态。
\textbf{15.} $169086.91 - 159855.97 = \qty{9230.94}{cm^{-1}}$：\qty{1083.3}{nm}，位于
近红外区。
\textbf{16.} $171134.89 - 166277.44 = \qty{4857.45}{cm^{-1}}$: \qty{2058.7}{nm}.
\textbf{17.} $\qty{171134.89}{cm^{-1}}$：\qty{58.43}{nm}，位于真空
紫外区（被空气吸收）。
\textbf{18.} 每一类都有自己的谱线和最低态，两类之间
不能通过光跃迁相互联系：它们的行为如同两种各有光谱的气体。
\textbf{19.} $E_\pm = E_{1s} + E_{2s} + J \pm K$（单重态取 $+$，三重态取 $-$）。
\textbf{20.} 三重态：电子相遇时其空间函数为零，
因此排斥作用较小。
\textbf{21.} \qty{6421.47}{cm^{-1}}, $6421.47/8065.54 = \qty{0.796}{eV}$.
\textbf{22.} 间隔 $\qty{2047.98}{cm^{-1}} = \qty{0.254}{eV}$，故 $K(1s,2p) =
\qty{0.127}{eV}$.
\textbf{23.} $2s$ 穿入 $1s$ 区域，与之重叠的程度
高于 $2p$。
\textbf{24.} 是的：两种构型中，三重态（$S$ 较大）都最低。
\textbf{25.} $K = \frac12 \times \qty{0.796}{eV}$：\textbf{氦的交换
积分 $K(1s,2s)$ 为 \qty{0.398}{eV}}，以单重态–三重态间隔的
一半度量。
\end{solution}
